% Contenido del póster: Morello (2025), Common Market Law Review. DOI 10.54648/cola2025039
% Elaborado con apoyo de IA a partir del texto completo del artículo, 2026-11-25.
% El estudiante debe leer el texto completo y reescribir con sus palabras.
\def\DOI{10.54648/cola2025039}
\def\FuenteSintesis{texto completo del artículo}
\def\TituloEs{Silos que se mezclan: el Reglamento de crowdfunding y la transformación del derecho financiero minorista de la Unión Europea}
\def\TituloOriginal{Blending silos. The Crowdfunding Regulation and the transformation of EU retail financial law}
\def\Autores{Morello}
\def\Anio{2025}
\def\Revista{Common Market Law Review}
\def\TituloGrande{\fontsize{58}{70}\selectfont}  % título largo: letra algo menor para que el póster quepa en una página

\def\Problema{En la Unión Europea, el crédito al consumidor y la inversión se regulan por separado, como «silos». El crowdfunding cruza esa frontera, porque combina crédito e inversión. Se critica que el Reglamento (UE) 2020/1503 deje fuera el crédito al consumidor, pero más allá de una lectura superficial se relaciona de cerca con el derecho del consumidor.}
\def\Objetivo{Mostrar hasta qué punto el Reglamento de crowdfunding acerca el derecho financiero al derecho del consumidor, y por qué ese acercamiento queda incompleto.}
\def\Metodo{Análisis jurídico que combina el derecho privado europeo y la regulación financiera. Primero describe la convergencia del derecho financiero minorista; luego examina las reglas que protegen al inversionista no sofisticado, y por último las compara con las del derecho del consumidor y del mercado de valores.}
\def\Resultados{\begin{itemize}
  \item El Reglamento cubre préstamos y valores de empresas, hasta 5 millones de euros en 12 meses, y excluye a los consumidores como receptores.
  \item El inversionista no sofisticado tiene una prueba de conocimientos, una simulación de pérdidas y cuatro días de reflexión.
  \item Esas reglas funcionan como derecho de contratos y pueden sumarse a las del consumidor.
  \item La protección es incompleta: un resultado desfavorable genera una advertencia, pero no impide invertir.
\end{itemize}}
\def\Conclusion{El Reglamento es el paso más avanzado de acercamiento entre el derecho del consumidor y la regulación financiera, pero no los une: uno busca equilibrar a las partes y la otra, la eficiencia del mercado.}
\def\Relevancia{La Unión Europea y el Perú regulan la misma actividad de forma distinta. El Reglamento europeo solo cubre proyectos de empresas; en el Perú, el reglamento de la SMV admite además proyectos personales, con topes de 50 UIT para ellos y de 500 UIT para los empresariales desde noviembre de 2025. Con datos de la SMV, la tasa de colocación fue de \cifra{colocacionEmp}~\% en los proyectos empresariales y de \cifra{colocacionPer}~\% en los personales. El estudio sugiere mirar no solo los topes, sino también qué protección recibe quien invierte.}
\def\PalabrasClave{crowdfunding, regulación, protección del inversionista, derecho del consumidor, Unión Europea}
% --- Ficha del estudio (herramientas/posteres/ficha_articulos.py) ---
\def\Ambito{Unión Europea (Reglamento 2020/1503)}
\def\Enfoque{Análisis jurídico: derecho financiero y derecho del consumidor}
\def\FiguraPeru{2}
% --- fin de la ficha ---
